% A integrar en la sección 5 junto con los modelos no lineales del rol 3.
\section{Metodología y modelos}
\subsection{Baseline de persistencia}
La referencia predice la siguiente hora con la medición actual: $\hat y_t=\mathrm{PM}_{2{,}5}(t)$. No requiere ajuste y representa una referencia exigente para una serie con dependencia temporal. Se evalúa sobre las mismas filas que los modelos entrenados.

\subsection{Modelos lineales regularizados}
Ridge y Lasso predicen mediante una combinación lineal de las características de la sección 4. Ridge penaliza la suma de cuadrados de los coeficientes, lo que reduce su magnitud; Lasso penaliza la suma de sus valores absolutos y puede llevar algunos exactamente a cero. En ambos casos $\alpha$ controla la intensidad de regularización. Se comparan versiones con nueve variables de historial y versiones que añaden el calendario. Las nueve variables de historial se estandarizan dentro del pipeline; la hora se representa con seno y coseno y el día con indicadores, con lunes como referencia. El ajuste temporal y la selección se describen en la sección 6.

En validación, la persistencia obtiene MAE 3,378, RMSE 5,143, MSE 26,446 y $R^2=0{,}700$. Ridge sin calendario logra MAE 3,141 y Lasso sin calendario 3,142; con calendario alcanzan 3,205 y 3,205, respectivamente. En esta comparación el calendario no mejora el MAE, por lo que no hay evidencia de aporte en este bloque para estas variantes. Ridge sin calendario, seleccionado por MAE de validación, mejora un 7,0~\% frente a persistencia en ese bloque. En prueba alcanza MAE 2,626 frente a 2,829 de persistencia (mejora de 7,2~\%); sus RMSE son 4,766 y 5,265, respectivamente. La comparación procede de la misma muestra y no implica que el resultado se mantenga en otros periodos.

El $\alpha$ seleccionado es 0,001 para ambas variantes de Ridge y 0,01 para ambas de Lasso. El valor de Ridge está en el límite inferior de la malla, por lo que no se descarta que una regularización menor mejore la validación temporal. En Lasso sin calendario ningún coeficiente es cero; con calendario se anulan los indicadores de martes, miércoles, jueves y sábado. El coeficiente de la medición actual es positivo y el de mayor magnitud entre las variables estandarizadas del historial; los coeficientes de estas variables correlacionadas deben interpretarse conjuntamente y no como efectos causales. Los indicadores de día no están estandarizados, por lo que sus magnitudes no se comparan directamente con las del historial. Las tablas completas de métricas, búsqueda y coeficientes se encuentran en \texttt{data/processed/lineales\_*.csv}.
